\SourceLeaf{394}%
woraus sich nach den entsprechenden Rechnungen leicht ergibt
\[
\begin{aligned}
&4\theta\sin(A''-A')\sin(A-A'')\sin(A'-A)\\
&\qquad=b\sin2(A''-A')+b'\sin2(A-A'')+b''\sin2(A'-A)
\end{aligned}
\]
\[
\begin{aligned}
&4uu\sin(A''-A')^2\sin(A-A'')^2\sin(A'-A)^2\\
&\qquad=bb\sin(A''-A')^2\\
&\qquad\quad+b'b'\sin(A-A'')^2\\
&\qquad\quad+b''b''\sin(A'-A)^2\\
&\qquad\quad+2b'b''\cos(A''-A')\sin(A-A'')\sin(A'-A)\\
&\qquad\quad+2bb''\sin(A''-A')\cos(A-A'')\sin(A'-A)\\
&\qquad\quad+2bb'\sin(A''-A')\sin(A-A'')\cos(A'-A)
\end{aligned}
\]
und daraus
\[
\begin{aligned}
&4(\theta\theta-uu)\sin(A''-A')^2\sin(A-A'')^2\sin(A'-A)^2\\
&\qquad=-bb\sin(A''-A')^4\\
&\qquad\quad-b'b'\sin(A-A'')^4\\
&\qquad\quad-b''b''\sin(A'-A)^4\\
&\qquad\quad+2b'b''\sin(A-A'')^2\sin(A'-A)^2\\
&\qquad\quad+2bb''\sin(A''-A')^2\sin(A'-A)^2\\
&\qquad\quad+2bb'\sin(A''-A')^2\sin(A''-A)^2
\end{aligned}
\]
\par
Ich habe diese Formeln hier aufgenommen, weil sie bisweilen auch in anderen Fällen nützlich sind. Man sieht leicht, dass der Ausdruck auf der rechten Seite das Produkt der vier Faktoren ist.
\[
\begin{array}{l}
+\BareRoot{b}.\sin(A''-A')+\BareRoot{b'}.\sin(A-A'')+\BareRoot{b''}.\sin(A'-A)\\
-\BareRoot{b}.\sin(A''-A')+\BareRoot{b'}.\sin(A-A'')+\BareRoot{b''}.\sin(A'-A)\\
+\BareRoot{b}.\sin(A''-A')-\BareRoot{b'}.\sin(A-A'')+\BareRoot{b''}.\sin(A'-A)\\
+\BareRoot{b}.\sin(A''-A')+\BareRoot{b'}.\sin(A-A'')-\BareRoot{b''}.\sin(A'-A)
\end{array}
\]
\par
Setzt man hier für \(b,b',b''\) ihre Werte ein und schreibt der Kürze halber
\[
a\sin(A''-A')+a'\sin(A-A'')+a''\sin(A'-A)=M
\]
so ist
\[
(\theta\theta-uu)\sin(A''-A')^2\sin(A-A'')^2\sin(A'-A)^2
\]
gleich dem Produkt der vier Faktoren
